\maketitle
\thispagestyle{plain}
\begin{abstract}

针对圆柱药材在时变烘房中的升温、失水和径向收缩过程，本文以干基含水率守恒为主线建立径向热湿输运模型。固定几何采用圆柱坐标，收缩几何通过材料坐标映射到固定参考域；烘房温度、环境水分量与实测半径由分段线性函数输入。空间离散采用环形有限体积法，时间推进采用BDF方法。

\textbf{问题一}采用常热物性和含水率相关扩散率。30 min时，中心与表面温度分别为$33.5753\ ^\circ\mathrm C$和$36.7856\ ^\circ\mathrm C$，中心含水率为$2.5500$，表面降至$1.5102\ \mathrm{kg/kg}$。热量已传入主体区域，而水分变化仍集中在表层。

\textbf{问题二}代入附录三变物性，建立由物性系数介导的双向热湿反馈模型。3 h时，中心与表面含水率分别为$1.7662$和$1.0081$；此时温场已接近环境平台，含水率梯度仍然明显，因而不能用温度稳定判断干燥完成。

\textbf{问题三}以全域最大含水率定义终止事件，得到等号临界时刻$57.4740$ h；首个严格低于阈值的整数秒为$206907$ s。若改用截面平均值或表面值判停，将分别提前约$21.45$ h和$44.92$ h。终点对扩散率尺度最敏感，后期干燥主要受内部输运制约。

\textbf{问题四}令$\xi=r/R(t)$随材料运动，将移动域转化为固定域。临界时刻为$51.0920$ h，此时半径为$1.2000$ cm。四组条件对照表明，问题三与问题四相差的$6.3820$ h同时受到物性和几何的影响；在给定外生半径轨迹下，两者还存在明显交互。

解析基线、网格加密、独立有限元和累计守恒检验均支持上述数值结果。问题三、四对扩散率尺度的无量纲灵敏度分别为$-0.8931$和$-0.8885$；相比数值离散，平衡含水率映射、长期环境和经验物性对实际预测的影响更值得关注。

\keywords{药材烘干；热湿耦合；有限体积法；移动边界；全域事件}
\end{abstract}
